% ------------------------------------------------------------------------
% file `revisions_juin-exercise-10-exercise-body.tex'
%   in folder `xsim/'
%
%     exercise of type `exercise' with id `10'
%
% generated by the `exercise' environment of the
%   `xsim' package v0.20c (2021/02/03)
% from source `revisions_juin' on 2023/06/12 on line 167
% ------------------------------------------------------------------------
	Vrai ou faux? Justifie dans tous les cas.
	\begin{enumerate}
		\item -3 est un entier
		\item L'addition admet 0 comme élément absorbant
		\item Dans un couple de coordonnées, le premier élément s'appelle l'ordonnée
		\item 5 et -5 ont la même valeur absolue
		\item Les diagonales d'un parallélogramme se coupent en leur milieu et sont isométriques.
		\item Il n'existe aucun triangle qui soit rectangle et isocèle
		\item 2,5 est un nombre naturel
		\item Pour calculer le coefficient de proportionnalité, on prend la première grandeur qu'on divise par la deuxième
		\item \(3x + 7 = 3x + 7\)
		\item Si j'ai 4 côtés et des diagonales isométriques, je suis un rectangle
		\item \(5 \cdot 2x \cdot 1 = 5 \cdot 2x\)
		\item La droite qui coupe un angle en deux parts égales est une médiane
		\item \((-2)^3=8\)
		\item L'opposé de 3 est 6
		\item Dans une translation, l'élément caractéristique est un vectorien
	\end{enumerate}
